\documentclass[11pt]{article}
\usepackage[a4paper,margin=25mm]{geometry}
\usepackage{amsmath,amssymb,bm,booktabs,longtable,hyperref}
\title{Corrected Two-Row Registry Energy and the Limits of an Ideal-Slip Extension}
\author{Donghyeok Kim \and Minsoo Kang \and Junhyeok Park}
\date{September 2026}
\newcommand{\dd}{\,\mathrm d}
\newcommand{\epsc}{\varepsilon_c}
\newcommand{\epsw}{\varepsilon_w}
\newcommand{\sigLJ}{\sigma_{\mathrm{LJ}}}
\begin{document}
\maketitle

\begin{abstract}
This note conservatively corrects the imported 23-page two-row slip-energy
derivation. The shifted Epstein--Hurwitz sum and its Poisson--Bessel
representation are retained. The coefficient and well-depth Lennard--Jones
conventions are separated, the energy is identified as cross-row energy per
repeat, and the unsupported addition of collinear bulk-chain and cross-row
energies is removed. An unwrapped registry index is retained only as a signed
translation counter: its change is not automatically an irreversible plastic
event. The construction remains an archived ideal-registry mechanism and is
not an active or quantitative plasticity law for aluminum.
\end{abstract}

\section{Scope and status}

The active repository theory remains one-dimensional and normal-only, with
state coordinate $a$ and density $P(a,t)$. The tangential registry $s$,
resolved shear $\tau$, and well index $z$ in this note are an inactive archive.
No active FEM element supplies shear stress to this model.

\section{Two distinct generalized Lennard--Jones conventions}

The active coefficient convention is
\begin{equation}
v(r)=\epsc\left[\left(\frac{\sigLJ}{r}\right)^m
-\left(\frac{\sigLJ}{r}\right)^n\right],\qquad m>n>1.
\end{equation}
Here $\epsc$ is a coefficient and is not the pair-well depth. If a separate
well-depth convention is desired, define
\begin{equation}
C_{m,n}=\frac{m}{m-n}\left(\frac{m}{n}\right)^{n/(m-n)},
\qquad \epsc=C_{m,n}\epsw.
\end{equation}
Then $v(r_e)=-\epsw$ at
$r_e=\sigLJ(m/n)^{1/(m-n)}$. The symbols $\epsc$ and $\epsw$ must never be
interchanged silently.

\section{Normal chain retained as a separate geometry}

For $N$ collinear atoms at spacing $a$,
\begin{equation}
E_N(a)=\sum_{k=1}^{N-1}(N-k)v(ka),
\qquad
U_\infty(a)=\epsc\left[\zeta(m)\left(\frac{\sigLJ}{a}\right)^m
-\zeta(n)\left(\frac{\sigLJ}{a}\right)^n\right].
\end{equation}
$U_\infty$ is energy per atom of a homogeneous collinear chain. It is not an
interface slip energy.

\section{Exact two-row cross energy per repeat}

Let the lower row contain $pb\bm d$, $p\in\mathbb Z$, and let one upper
reference atom be separated by $a\bm n+s\bm d$, with
$\bm d\cdot\bm n=0$. The pair distances are
\begin{equation}
r_p(a,s)=\sqrt{a^2+(pb+s)^2}.
\end{equation}
The cross-row energy per upper atom, equivalently per commensurate row repeat,
is
\begin{equation}
W(a,s)=\epsc\sum_{p\in\mathbb Z}\left[
\frac{\sigLJ^m}{[a^2+(pb+s)^2]^{m/2}}
-\frac{\sigLJ^n}{[a^2+(pb+s)^2]^{n/2}}
\right].
\end{equation}
It is not the divergent total energy of two infinite rows. Set
$\delta=s/b$, $\eta=a/b$ and
\begin{equation}
\mathcal Z_\nu(\delta,\eta)=\sum_{p\in\mathbb Z}
[(p+\delta)^2+\eta^2]^{-\nu}.
\end{equation}
For $\mathrm{Re}\,\nu>1/2$,
\begin{equation}
\begin{aligned}
\mathcal Z_\nu={\sqrt\pi\over\Gamma(\nu)}\Bigg[
&\Gamma\left(\nu-\frac12\right)\eta^{1-2\nu}\\
&+4\sum_{\ell=1}^\infty\cos(2\pi\ell\delta)
\left(\frac{\pi\ell}{\eta}\right)^{\nu-1/2}
K_{\nu-1/2}(2\pi\ell\eta)\Bigg].
\end{aligned}
\end{equation}
Consequently
\begin{equation}
W(a,s)=\epsc\left[
\left(\frac{\sigLJ}{b}\right)^m\mathcal Z_{m/2}
-\left(\frac{\sigLJ}{b}\right)^n\mathcal Z_{n/2}\right],
\qquad W(a,s+b)=W(a,s).
\end{equation}
This is an exact identity for the stated ideal two-row pair geometry, not for
an FCC half-space.

\section{Reference registry and traction}

A reference $s_0$ may be called a perfect registry only after checking
\begin{equation}
\partial_sW(a,s_0)=0,\qquad \partial_s^2W(a,s_0)>0
\end{equation}
at the specified $a$. Otherwise $W(a,s)-W(a,s_0)$ is merely a signed energy
difference. If a physical interface area $A_{\rm rep}$ per row repeat has been
defined crystallographically, then
\begin{equation}
\gamma(a,s)={W(a,s)-W(a,s_0)\over A_{\rm rep}},
\qquad \tau_{\rm int}=\partial_s\gamma.
\end{equation}
A one-dimensional row does not determine $A_{\rm rep}$ by itself.

\section{No unsupported mixed patch energy}

The collinear $U_\infty(a)$ and cross-row $W(a,s)$ describe different pair
geometries. Therefore the expression
$N_AU_\infty+A_0\gamma$ is not adopted. For the toy two-row interface alone,
one may define, under an explicit count $N_{\rm rep}$,
\begin{equation}
\Phi_{\rm row}=N_{\rm rep}[W(a,s)-W(a_0,s_0)]
-\sigma A_0(a-a_0)-\tau A_0(s-s_0).
\end{equation}
This remains a representative construction. A quantitative coupled model
must instead start from one FCC half-space Hamiltonian and perform a single,
explicit bulk/interface subtraction.

\section{Unwrapped registry is not automatic irreversibility}

Write the kinematic identity
\begin{equation}
s=zb+\widetilde s,qquad z\in\mathbb Z,qquad 0\le\widetilde s<b.
\end{equation}
$z$ records signed net translations. A conservative periodic landscape permits
both $z\to z+1$ and $z\to z-1$. Thus a jump is not automatically permanent.
For a declared homogenization thickness $h_{\rm slip}$, a candidate residual
shear strain is
\begin{equation}
\gamma_p={b\over h_{\rm slip}}\langle z\rangle,
\end{equation}
but it is plastic only if the $z$ distribution remains shifted after complete
unloading and the specified relaxation observation time.

A thermal master equation must include both directions,
\begin{equation}
\partial_tp_z=-\partial_aJ_{a,z}
+k_{z-1,+}p_{z-1}+k_{z+1,-}p_{z+1}
-(k_{z,+}+k_{z,-})p_z,
\end{equation}
with rates satisfying local detailed balance for the stated bath and tilted
energy. A single periodic potential supplies neither dislocation hardening nor
a guarantee of net slip under symmetric cyclic loading.

\section{Probability and boundary corrections}

For an absorbing intact density $\rho$, define raw moments and survival by
\begin{equation}
S=\iint\rho\dd a\dd s,qquad M_k=\iint a^k\rho\dd a\dd s.
\end{equation}
The surviving-population mean and variance are
\begin{equation}
\langle a\rangle_c={M_1\over S},
\qquad \operatorname{Var}_c[a]={M_2\over S}-\left({M_1\over S}\right)^2.
\end{equation}
An absorbing boundary has no intact mass outside it. Hence initiation is
measured by outgoing flux, $S$, hazard, and $1-S$, not by integrating the
absorbed density beyond the boundary. The active normal model's separately
declared $a_c$ first-passage definition is not changed by this archive.

\section{Dynamics and energy destination}

A projected atomistic model generally has a potential of mean force and a
matrix memory kernel. Replacing these by the displayed energy, diagonal
constant friction, and white noise is an assumption requiring MD validation.
The overdamped limit requires momentum and memory relaxation to be fast not
only relative to loading but also relative to resolved configurational motion.

With Langevin friction, positive loop work is transferred to unresolved
atomic and phonon modes. Without that bath, fracture, or another irreversible
state change, the fixed multiwell energy is conservative and cannot by itself
produce steady dissipative plastic hysteresis.

\section{Requirements for aluminum activation}

The two-row result is not a quantitative aluminum law. Activation requires a
declared FCC plane and direction, a complete validated Al EAM potential or DFT
generalized-stacking-fault surface, the corresponding representative area,
MD-derived memory/mobility, residual-slip validation, and additional physics
for dislocation storage and cyclic hardening. Until then this note remains an
ideal single-registry mechanism archive.

\begin{thebibliography}{9}
\bibitem{dlmf} NIST Digital Library of Mathematical Functions,
Section 10.32, Integral Representations of Modified Bessel Functions.
\bibitem{dawbaskes} M. S. Daw and M. I. Baskes, ``Embedded-atom method:
Derivation and application to impurities, surfaces, and other defects in
metals,'' \emph{Phys. Rev. B} 29, 6443--6453 (1984).
\bibitem{lu} G. Lu, N. Kioussis, V. V. Bulatov, and E. Kaxiras,
``Generalized-stacking-fault energy surface and dislocation properties of
aluminum,'' \emph{Phys. Rev. B} 62, 3099--3108 (2000).
\end{thebibliography}
\end{document}
